
The average S/N for a given PTA and a predicted  SGWB signal following the method in~\cite{SiemensEtAl:2013}. 
The power spectrum for a SGWB is given by
\begin{equation}
\Pg (f) = \frac{A^2}{24 \pi^2} \left(\frac{f}{\fref}\right)^{2\alpha} f^{-3}\,,
\label{eqn:Pgw}
% equation 16 in SiemansEtAl:2013
\end{equation}
where $f$ is the gravitational-wave frequency, $\fref$ is a reference frequency chosen as $1\,{\rm yr}^{-1}$, $A$ is the amplitude of the background at $f=\fref$, and $\alpha$ is the slope of the spectrum. 
Throughout this work we assume $a=-2/3$ consistent with a stochastic background from circular massive black hole binaries and $A = 2 \times 10^{-15}$ consistent with recent speculative hints of detection~\citep{GoncharovEtAlPPTA:2021,NANOGrav12.5GWB:2020}.

The average SNR for a PTA with total timespan $T$ is given by
\begin{equation}
\left< \rho \right> = \left( 2 T \sum_{IJ} \chiIJhd^2 
                             \int_{f_L}^{f_H} {\rm d}f  
                             \frac{\Pg^2 (f)}{\PI(f)\PJ(f)}
                            \right)^{1/2} \,,
\label{eqn:snrPIPJPg}
\end{equation}
where $\sum_{IJ}$ is the sum over all pulsar pairs ($I \neq J$) and $P_I$ is the power spectrum of the $I{\rm th}$ pulsar. 
We assume that the timespan of observations, $T$ is the same for all pulsars. 
The integration limits are given by $f_L=1/T$ and $f_H=c/2$ where $c$ is the cadence of observations, which we also assume to be the identical for all pulsars in the array. 
The Hellings and Downs coefficients $\chi_{IJ}$ describe the expected angular correlation produced by the SGWB signal in pairs of pulsars and is given by~\citep{HellingsDowns:1983}
\begin{eqnarray}
\chiIJhd &=& \frac{1}{2} 
           - \frac{1}{4} \left(\frac{1 - \cos(\zeta)}{2}\right) \nonumber\\
           & &+ \frac{3}{2} \left(\frac{1 - \cos(\zeta)}{2}\right) 
           \log \left( \frac{1 - \cos(\zeta)}{2} \right)\,,
\label{eqn:hd}
\end{eqnarray}
for pulsars separated by angle $\zeta$.

If we assume evenly sampled data and that there is only white noise in the pulsars then $\PI$ is the sum of Eqn.~\ref{eqn:Pgw} plus white noise, 
\begin{equation}
\PI (f) = \frac{A^2}{24 \pi^2} \left(\frac{f}{\fref}\right)^{2\alpha} f^{-3} 
          + 2 \sigma_I^2 \Delta t\,,
\end{equation}
where $c=1/\Delta t$ and $\sigma_I$ is the timing precision for the $I{\rm th}$ pulsar. 

Pulsars may also contain spin noise (or red noise) and jitter noise due to changes in the pulse profile shape. 
The power spectral density of spin noise for the $I{\rm th}$ pulsar is given by 
\begin{equation}
% from Goncharov 2021 eqn 5
\PsnI (f) = \frac{\ARedI^2}{12 \pi^2} 
            \left( \frac{f}{\fref} \right)^{-\gammaRedI} {\rm yr}^3\,,
\end{equation}
where $\ARedI$ and $\gammaRedI$ are the amplitude and slope of the power law for the $I$th pulsar and ${\rm yr}$ is the number of seconds in a year. 
Jitter noise $\sigma_{\rm jit}$ is added in quadrature to the white noise part the the spectrum.
The resulting power spectrum is
\begin{eqnarray}
\PI (f) &=& \frac{A^2}{24 \pi^2} \left(\frac{f}{\fref}\right)^{2\alpha} f^{-3} \nonumber\\
          & &~+~ 2 \left( \sigma_I^2 + \sigma_{{\rm jit}, I}^2 \right) \Delta t \nonumber\\
          & &~+~ \frac{\ARedI^2}{12 \pi^2} 
          \left( \frac{f}{\fref} \right)^{-\gammaRedI} {\rm yr}^3\,,
\label{eqn:PIGWWR}
\end{eqnarray}
which includes SGWB, white noise, spin noise, and jitter noise.
Equations~\ref{eqn:snrPIPJPg} and~\ref{eqn:PIGWWR} can therefore be used to compute the average PTA SNR for a given observation strategy and assumed SGWB signal. 
\begin{comment}
\begin{equation}
\left< \rho \right> = \left( 2 T \sum_{IJ} \chi_{IJ}^2
                             \int_{f_L}^{f_H} {\rm d}f  
                             \frac{\Pg^2 (f)}
                                  {[\Pg (f) + 2\sigma_I^2\Delta t + \Psn (f)]
                                   [\Pg (f) + 2\sigma_J^2\Delta t + \Psn (f)]} 
                            \right)^{1/2} \,.
\label{eqn:snrWhiteNoiseOnly}
\end{equation}
\end{comment}

